France and the Netherlands include electoral units located outside mainland
Europe.  These are represented under their respective country codes
(\texttt{FR}, \texttt{NL}) using dedicated NUTS-like parent codes that extend
the standard GISCO hierarchy.

\paragraph{France.}
French overseas units are grouped under two extended NUTS~1 branches.
The \texttt{FRY} branch covers the five French overseas departments and regions
(\textit{DOM-ROM}), which are full NUTS~2 and NUTS~3 regions in the standard
Eurostat classification; their municipalities are stored as ordinary LAU entries
in \bluepath{geo/communes.csv}.
The \texttt{FRZ} branch covers the French overseas collectivities
(\textit{collectivités d'outre-mer}), which have no official NUTS codes.
BLUE\_DB introduces a parallel NUTS-like hierarchy (\texttt{FRZ1}--\texttt{FRZ5})
to allow territory-level aggregation; the municipalities in this branch are
treated as \texttt{Overseas Country or Territory} entries in
\bluepath{geo/special.csv}, with their \bluevar{parent} fields pointing into
the extended hierarchy.

\paragraph{Netherlands.}
The three Caribbean special municipalities (Bonaire, Sint Eustatius, Saba ---
the \textit{BES islands}) are assigned the parent codes \texttt{NL901},
\texttt{NL902}, and \texttt{NL903} respectively, with the  NUTS~2
pseudo-region \texttt{NL90} representing Caribbean Netherlands (\emph{Caribisch Nederland}).

The following table lists the non-standard NUTS-like parent codes introduced
by BLUE\_DB for the \texttt{FRZ} collectivités and the Dutch BES islands,
with the number of municipalities referenced under each.
